% ============================================================
\setpartpreamble{\partintro
  {Generating schema changes, extending the dotnet ef tooling and supporting compiled
  models.}
  {
    \item Chapter~\ref{ch:how-migrations-work}: How migrations work
    \item Chapter~\ref{ch:declarative-sql-from-annotations}: Declarative SQL from
    annotations
    \item Chapter~\ref{ch:custom-migration-operations}: Custom migration operations
    \item Chapter~\ref{ch:design-time-services}: Design-time services
    \item Chapter~\ref{ch:compiled-models-trimming-and-nativeaot}: Compiled models, trimming
    and NativeAOT
    \item Chapter~\ref{ch:scaffolding-reverse-engineering}: Scaffolding (reverse
    engineering)
  }}
\part{Migrations and design time}\label{part:migrations}

% ------------------------------------------------------------
\chapter{How migrations work}\label{ch:how-migrations-work}
\versionrisk{stable}

Migrations compare two models and produce operations and SQL from them. A plugin can hook in
at several points; to choose the right one you need to know the flow:

\begin{flowdiagram}
  \flowstep{Design-time model}{Model including all design-time annotations}
  \flowstep{IRelationalAnnotationProvider}{Transfers annotations to tables, columns and
  indexes of the relational model}
  \flowstep{IMigrationsModelDiffer}{Compares the old (snapshot) and new relational model}
  \flowstep{MigrationOperations}{CreateTable, AlterTable, ... with .Annotation /
  .OldAnnotation}
  \flowstep{IMigrationsCodeGenerator}{Writes migration, designer file and snapshot (dotnet
  ef)}
  \flowstep{IMigrationsSqlGenerator}{Produces SQL when applying or scripting}
\end{flowdiagram}

There are two basic strategies for plugins. \textbf{Declarative}: the plugin describes the
desired state as annotations; EF detects changes by itself, and the plugin only has to
produce SQL. \textbf{Imperative}: the plugin offers its own operations that users write into
migrations by hand. Good plugins use both: declarative for everything that belongs to the
model, imperative for everything else.

\begin{note}
  \code{Database.\allowbreak{}Ensure\allowbreak{}Created(\allowbreak{})} uses the same
  pipeline: it computes the diff from the empty state to the design-time model and executes
  the SQL of the migrations generator. A plugin's custom migration SQL therefore takes
  effect there too, which simplifies integration tests considerably
  (chapter~\ref{ch:test-strategy}).
\end{note}

% ------------------------------------------------------------
\chapter{Declarative SQL from annotations}\label{ch:declarative-sql-from-annotations}
\versionrisk{volatile}

\section{The annotation provider}

The differ compares the \emph{relational} model (tables, columns), not the entity types.
Which annotations a table carries is decided by the
\code{IRelational\allowbreak{}Annotation\allowbreak{}Provider}. There is only one (a
singleton per internal service provider), and it belongs to the provider; the plugin therefore derives from the Npgsql implementation:

% source: src/Acme.EntityFrameworkCore.Tenancy.PostgreSQL/Metadata/Internal/TenancyAnnotationProvider.cs
\begin{sourcecode}[Metadata/\allowbreak{}Internal/\allowbreak{}Tenancy\allowbreak{}Annotation\allowbreak{}Provider.\allowbreak{}cs]{csharp}
public class TenancyAnnotationProvider : NpgsqlAnnotationProvider
{
    public TenancyAnnotationProvider(RelationalAnnotationProviderDependencies dependencies)
        : base(dependencies)
    {
    }

    public override IEnumerable<IAnnotation> For(ITable table, bool designTime)
    {
        foreach (var annotation in base.For(table, designTime))
            yield return annotation;

        if (!designTime)          // migrations only: the runtime model stays lean
            yield break;

        // TPH/table splitting: the type with RLS configuration, not simply the first
        var entityType = table.EntityTypeMappings
            .Select(m => (IEntityType)m.TypeBase)
            .FirstOrDefault(e => e.HasRlsConfiguration())
            ?? (IEntityType)table.EntityTypeMappings.First().TypeBase;
        foreach (var annotation in entityType.GetAnnotations())
        {
            if (!annotation.Name.StartsWith(RlsAnnotationNames.Prefix,
                                            StringComparison.Ordinal))
                continue;
            // Opt-out marker and explicit force: false stay in the model (they beat the
            // conventions), but "false" and "missing" must be the same table
            if (annotation.Name is RlsAnnotationNames.Enabled or RlsAnnotationNames.Forced
                && annotation.Value is false)
                continue;
            yield return annotation;
        }
    }
}
\end{sourcecode}

From here EF takes over: when a policy changes, the differ produces an
\code{Alter\allowbreak{}Table\allowbreak{}Operation} with new and old annotations. In the
generated migration it looks like this, without any custom code generator:

\begin{sourcecode}[Generated migration (\allowbreak{}excerpt)]{csharp}
migrationBuilder.AlterTable(name: "Orders")
    .Annotation("Acme:Rls:Policy:tenant_isolation", "{...new...}")
    .OldAnnotation("Acme:Rls:Policy:tenant_isolation", "{...old...}");
\end{sourcecode}

Because the annotation provider only passes the annotations on for \code{design\allowbreak{}Time},
they are missing from the runtime model that \code{context.\allowbreak{}Model} returns.

\begin{sourcecode}[Startup check of a plugin (\allowbreak{}WRONG)]{csharp}
// Does every isolated table have RLS? Read from the runtime model
var table = context.Model.GetRelationalModel().FindTable("Orders", null);
var enabled = table?[RlsAnnotationNames.Enabled];          // always null
\end{sourcecode}

\begin{failure}[What happens if plugin code reads the runtime model instead of the design-time model?]
  \textbf{Observable.} Code that reads the plugin's schema annotations from the relational
  runtime model finds none: a check reports nothing, a self-generated script contains no
  policies. EF protects the differ itself: \code{Get\allowbreak{}Differences} with the
  runtime model fails with a hint to \code{IDesign\allowbreak{}Time\allowbreak{}Model}.

  \textbf{Cause.} The runtime model is lean. The annotation provider passes annotations for
  tables on only at design time, so that they do not burden the run time.

  \textbf{Solution.} For everything related to the schema, use
  \code{context.\allowbreak{}Get\allowbreak{}Service<\allowbreak{}IDesign\allowbreak{}Time\allowbreak{}Model>(\allowbreak{}).\allowbreak{}Model}.

  \textbf{Test:}
  \code{Failure\allowbreak{}Case\allowbreak{}Tests.\allowbreak{}Table\allowbreak{}\_annotations\allowbreak{}\_of\allowbreak{}\_the\allowbreak{}\_runtime\allowbreak{}\_model\allowbreak{}\_…}
  and \code{Differ\allowbreak{}\_rejects\allowbreak{}\_the\allowbreak{}\_runtime\allowbreak{}\_model\allowbreak{}\_…}.
\end{failure}

\section{The SQL generator}

There is also only one \code{IMigrations\allowbreak{}Sql\allowbreak{}Generator}. The plugin
derives from the Npgsql generator and extends \code{Create\allowbreak{}Table} and
\code{Alter\allowbreak{}Table}:

% source: src/Acme.EntityFrameworkCore.Tenancy.PostgreSQL/Migrations/Internal/TenancyMigrationsSqlGenerator.cs (strings)
\begin{sourcecode}[Migrations/\allowbreak{}Internal/\allowbreak{}Tenancy\allowbreak{}Migrations\allowbreak{}Sql\allowbreak{}Generator.\allowbreak{}cs (\allowbreak{}excerpt)]{csharp}
protected override void Generate(
    CreateTableOperation operation, IModel? model,
    MigrationCommandListBuilder builder, bool terminate = true)
{
    // Unterminated SQL cannot take further statements (Npgsql handles comments the same way)
    if (!terminate && operation.HasRlsConfiguration())
        throw new ArgumentException(
            "RLS configuration cannot be emitted in an unterminated CREATE TABLE.");

    base.Generate(operation, model, builder, terminate);
    if (terminate)
        GenerateRowLevelSecurity(operation.Name, operation.Schema, previous: null,
            operation, builder);
}

protected override void Generate(
    AlterTableOperation operation, IModel? model, MigrationCommandListBuilder builder)
{
    base.Generate(operation, model, builder);
    if (_moved.Contains(operation))
        GenerateRowLevelSecurityState(Delimit(operation.Name, operation.Schema),
            operation.OldTable, operation, builder);
    else
        GenerateRowLevelSecurity(operation.Name, operation.Schema, operation.OldTable,
            operation, builder);
}
\end{sourcecode}

\begin{sourcecode}[Migrations/\allowbreak{}Internal/\allowbreak{}Tenancy\allowbreak{}Migrations\allowbreak{}Sql\allowbreak{}Generator.\allowbreak{}cs (\allowbreak{}shortened, simplified)]{csharp}
protected virtual void GenerateRowLevelSecurity(string tableName, string? schema,
    IReadOnlyAnnotatable? previous, IReadOnlyAnnotatable current,
    MigrationCommandListBuilder builder)
{
    var table = Delimit(tableName, schema);
    var oldPolicies = previous?.GetRlsPolicies() ?? new Dictionary<string, string>();
    var newPolicies = current.GetRlsPolicies();

    // 1. Drop removed or changed policies
    foreach (var (name, json) in oldPolicies)
        if (!newPolicies.TryGetValue(name, out var n) || n != json)
            Statement(builder, $"DROP POLICY {Delimit(name)} ON {table}");

    // 2. RLS state (ENABLE/DISABLE, FORCE/NO FORCE)
    if ((previous?.IsRlsEnabled() ?? false) != current.IsRlsEnabled())
        Statement(builder, $"ALTER TABLE {table} "
            + (current.IsRlsEnabled() ? "ENABLE" : "DISABLE") + " ROW LEVEL SECURITY");
    // ... FORCE analogous

    // 3. Create new or changed policies
    foreach (var (name, json) in newPolicies)
        if (!oldPolicies.TryGetValue(name, out var o) || o != json)
            Statement(builder, CreatePolicySql(table, RlsPolicy.Deserialize(json)));
}

protected virtual void Statement(MigrationCommandListBuilder builder, string sql)
{
    builder.Append(sql).AppendLine(Dependencies.SqlGenerationHelper.StatementTerminator);
    EndStatement(builder);
}
\end{sourcecode}

\begin{datatable}{H{0.60}H{1.40}}
  \textbf{Detail} & \textbf{Reason} \\
  \tablehead
  DROP + CREATE instead of \code{ALTER POLICY} & \code{ALTER POLICY} can change neither
  \code{FOR} nor PERMISSIVE/RESTRICTIVE. DDL is transactional in PostgreSQL; no visible gap
  occurs. \\
  \tablerule
  \code{End\allowbreak{}Statement(\allowbreak{}builder)} & Ends the command like the base
  class does, honoring transaction suppression. \\
  \tablerule
  Own \code{Delimit} helpers & Npgsql's \code{Delimit\allowbreak{}Identifier} helpers are
  private. Always escape identifiers via
  \code{Dependencies.\allowbreak{}Sql\allowbreak{}Generation\allowbreak{}Helper}, never by
  string concatenation. \\
  \tablerule
  DropTable, RenameTable & Nothing to do: PostgreSQL removes or moves policies together with
  the table. \\
  \tablerule
  \code{\_moved} & \code{Alter\allowbreak{}Table\allowbreak{}Operations} whose policy
  statements were moved into operations of their own (next section); for them the generator
  only emits ENABLE/DISABLE and FORCE. \\
\end{datatable}

\subsection*{Policies and column changes of the same table}

A policy depends on the columns it names: \code{CREATE POLICY} needs them, and
\code{DROP COLUMN} fails while a policy uses them. EF's differ, however, orders a table's
\code{Alter\allowbreak{}Table\allowbreak{}Operation} before its new columns and after its dropped
ones. When an existing table gets a tenant column and the isolation in one migration, it yields
\code{Alter\allowbreak{}Table}, \code{Add\allowbreak{}Column}, \code{Create\allowbreak{}Index};
the reverse yields \code{Drop\allowbreak{}Index}, \code{Drop\allowbreak{}Column},
\code{Alter\allowbreak{}Table}. A generator that creates the policy right at the
\code{Alter\allowbreak{}Table\allowbreak{}Operation} fails in both directions: PostgreSQL reports
\code{column "TenantId" does not exist} or \code{cannot drop column TenantId \ldots{} because
other objects depend on it}.

The generator therefore also overrides \code{Generate} for the whole operation list.
\code{Move\allowbreak{}Policy\allowbreak{}Changes} replaces the policy statements of every changed
table with internal operations of their own: changed and removed policies are dropped before the
table's first column change, new and changed ones are created after its last. For the migration
above this gives:

\begin{sourcecode}[Generated SQL (\allowbreak{}shortened)]{sql}
ALTER TABLE "Products" ENABLE ROW LEVEL SECURITY;
ALTER TABLE "Products" FORCE ROW LEVEL SECURITY;
ALTER TABLE "Products" ADD "TenantId" uuid NOT NULL DEFAULT '00000000-...';
CREATE POLICY tenant_isolation ON "Products" AS PERMISSIVE FOR ALL USING (...) ...;
CREATE INDEX "IX_Products_TenantId" ON "Products" ("TenantId");
\end{sourcecode}

A renamed table keeps its identity; every statement uses the name that is valid at its position.
All statements stay in the migration's transaction: if a later step fails, the rollback restores
policy, RLS and column together, and the table is not left unprotected. The SQL Server package
solves the same question more coarsely: it drops the policies of all affected tables before the
first operation and recreates them after the last, because there \code{sp\_rename} and
\code{ALTER COLUMN} already fail on the schema-bound policy. PostgreSQL moves policies along with
renames itself; here the smaller window around the column changes is enough.

Two limits remain. A column change that does not change the policy text, for example another
type for the tenant column, produces no \code{Alter\allowbreak{}Table\allowbreak{}Operation};
PostgreSQL refuses \code{ALTER COLUMN \ldots{} TYPE} for a column a policy names, and the policy
has to be replaced by hand in the migration. And a custom operation with
\code{suppress\allowbreak{}Transaction: true} between the column changes of a table splits the
transaction; the rollback then only protects up to that operation.

% ------------------------------------------------------------
\chapter{Custom migration operations}\label{ch:custom-migration-operations}
\versionrisk{medium}

Some things do not belong in the model but in migrations: roles, grants, extensions, data
backfills. For these the plugin offers its own operations. They consist of three parts: the
operation class, an extension method on \code{Migration\allowbreak{}Builder} and the SQL
generation in the generator.

% source: src/Acme.EntityFrameworkCore.Tenancy.PostgreSQL/Migrations/Operations/TenancyOperations.cs
\begin{sourcecode}[Migrations/\allowbreak{}Operations/\allowbreak{}Tenancy\allowbreak{}Operations.\allowbreak{}cs (\allowbreak{}excerpt)]{csharp}
public class CreateApplicationRoleOperation : MigrationOperation
{
    public required string Name { get; init; }
    public bool CanLogin { get; init; }       // passwords do NOT belong in migrations
}
\end{sourcecode}

% source: src/Acme.EntityFrameworkCore.Tenancy.PostgreSQL/Migrations/Operations/TenancyMigrationBuilderExtensions.cs (strings)
\begin{sourcecode}[Migrations/\allowbreak{}Operations/\allowbreak{}Tenancy\allowbreak{}Migration\allowbreak{}Builder\allowbreak{}Extensions.\allowbreak{}cs (\allowbreak{}excerpt)]{csharp}
namespace Microsoft.EntityFrameworkCore.Migrations;      // always imported in migrations

public static class TenancyMigrationBuilderExtensions
{
    public static OperationBuilder<CreateApplicationRoleOperation> CreateApplicationRole(
        this MigrationBuilder migrationBuilder, string name, bool canLogin = false)
    {
        var operation = new CreateApplicationRoleOperation
        {
            Name = TenancyIdentifiers.ValidateIdentifier(name, "role name"),
            CanLogin = canLogin,
        };
        migrationBuilder.Operations.Add(operation);
        return new OperationBuilder<CreateApplicationRoleOperation>(operation);
    }
}
\end{sourcecode}

The generator dispatches operations through the central
\code{Generate(\allowbreak{}Migration\allowbreak{}Operation,\allowbreak{}
.\allowbreak{}.\allowbreak{}.\allowbreak{})} method. Unknown types go to the base class:

\begin{sourcecode}[Migrations/\allowbreak{}Internal/\allowbreak{}Tenancy\allowbreak{}Migrations\allowbreak{}Sql\allowbreak{}Generator.\allowbreak{}cs (\allowbreak{}excerpt, simplified)]{csharp}
protected override void Generate(
    MigrationOperation operation, IModel? model, MigrationCommandListBuilder builder)
{
    switch (operation)
    {
        case CreateApplicationRoleOperation createRole: Generate(createRole, builder); break;
        case DropApplicationRoleOperation dropRole:     Generate(dropRole, builder); break;
        case GrantTablePrivilegesOperation grant:       Generate(grant, builder); break;
        default: base.Generate(operation, model, builder); break;
    }
}

protected virtual void Generate(CreateApplicationRoleOperation operation,
    MigrationCommandListBuilder builder)
{
    var literal = Literal(operation.Name);           // 'shop_app'  (for the comparison)
    var identifier = Delimit(operation.Name);    // shop_app (CREATE ROLE, "..." if needed)
    var login = operation.CanLogin ? "LOGIN" : "NOLOGIN";

    builder
        .AppendLine("DO $$")
        .AppendLine("BEGIN")
        .AppendLine("    IF NOT EXISTS (SELECT FROM pg_catalog.pg_roles")
        .AppendLine($"                   WHERE rolname = {literal}) THEN")
        .AppendLine($"        CREATE ROLE {identifier} {login} NOSUPERUSER NOBYPASSRLS;")
        .AppendLine("    END IF;")
        .AppendLine("END")
        .Append("$$");
    builder.AppendLine(Dependencies.SqlGenerationHelper.StatementTerminator);
    EndStatement(builder);
}
\end{sourcecode}

\begin{important}[Literal or identifier]
  The role name appears twice, but in two forms: as a \textbf{string literal} in the
  comparison (\code{'shop\_\allowbreak{}app'}, produced by the string type mapping) and as
  an \textbf{identifier} in \code{CREATE ROLE} (\code{shop\_\allowbreak{}app}, in double
  quotes if it needs them). Mixing them up or building them by string concatenation produces
  invalid SQL or an injection hole. The extension method additionally validates the name.
\end{important}

\begin{tip}
  Write operations idempotently (\code{IF NOT EXISTS}, \code{DROP
  .\allowbreak{}.\allowbreak{}.\allowbreak{} IF EXISTS}). Roles exist cluster-wide and
  outlive databases; a migration that fails on its second run against a new database is a
  common annoyance.
\end{tip}

In a migration, the operations are used like EF's own:

\begin{sourcecode}[Migrations/\allowbreak{}20261021\_\allowbreak{}Initial.\allowbreak{}cs (\allowbreak{}hand-edited)]{csharp}
protected override void Up(MigrationBuilder migrationBuilder)
{
    // BEFORE the tables (chapter "Advanced migrations")
    migrationBuilder.CreateApplicationRole("shop_support");
    // ... operations generated by EF (tables, RLS via annotations) ...
    migrationBuilder.CreateApplicationRole("shop_app", canLogin: true);
    migrationBuilder.GrantTablePrivileges("shop_app");
}

protected override void Down(MigrationBuilder migrationBuilder)
{
    migrationBuilder.DropApplicationRole("shop_app");
    // ...
}
\end{sourcecode}

\begin{note}[When the differ should produce custom operations]
  If the \emph{differ} should produce custom operations automatically, you need a derived
  \code{IMigrations\allowbreak{}Model\allowbreak{}Differ} and an
  \code{ICSharp\allowbreak{}Migration\allowbreak{}Operation\allowbreak{}Generator} that
  writes the operation as C\#. Both are single services with a lot of internal API. The
  declarative route via annotations (chapter~\ref{ch:declarative-sql-from-annotations}) is
  almost always the better choice.
\end{note}

% ------------------------------------------------------------
\chapter{Design-time services}\label{ch:design-time-services}
\versionrisk{volatile}

The \code{dotnet ef} tools (adding migrations, writing snapshots, scaffolding, compiled
models) use their own service provider with design-time services. A plugin extends it
through \code{IDesign\allowbreak{}Time\allowbreak{}Services}.

\section{Loaded through the startup project}

An assembly attribute points to the class. The tooling, however, reads it only from the
startup project and from the migrations assembly, not from referenced packages; an attribute
in the plugin assembly itself would have no effect. The plugin therefore documents one line
that the application adds once to its startup project. The second argument restricts loading
to contexts with the matching provider. The consumer test
(page~\pageref{sec:consumer-test}) does exactly that:

% source: eng/consumer-test/template/DesignTimeServices.cs
\begin{sourcecode}[eng/\allowbreak{}consumer-test/\allowbreak{}template/\allowbreak{}Design\allowbreak{}Time\allowbreak{}Services.\allowbreak{}cs]{csharp}
[assembly: DesignTimeServicesReference(
    "Acme.EntityFrameworkCore.Tenancy.PostgreSQL.Design.TenancyDesignTimeServices, "
        + "Acme.EntityFrameworkCore.Tenancy.PostgreSQL",
    "Npgsql.EntityFrameworkCore.PostgreSQL")]
\end{sourcecode}

% source: src/Acme.EntityFrameworkCore.Tenancy.PostgreSQL/Design/TenancyDesignTimeServices.cs
\begin{sourcecode}[Design/\allowbreak{}Tenancy\allowbreak{}Design\allowbreak{}Time\allowbreak{}Services.\allowbreak{}cs (\allowbreak{}excerpt)]{csharp}
public class TenancyDesignTimeServices : IDesignTimeServices
{
    public void ConfigureDesignTimeServices(IServiceCollection services)
    {
        // Remove all registrations (EF default, Npgsql), then add ours:
        // independent of the order in which the tooling loads services
        services.RemoveAll<IAnnotationCodeGenerator>();
        services.AddSingleton<IAnnotationCodeGenerator, TenancyAnnotationCodeGenerator>();

        // ... likewise ICSharpRuntimeAnnotationCodeGenerator (compiled models); the
        // IDatabaseModelFactory (scaffolding), however, scoped as in Npgsql
    }
}
\end{sourcecode}

\section{Readable snapshots: the annotation code generator}

By default EF writes unknown annotations as
\code{Has\allowbreak{}Annotation(\allowbreak{}.\allowbreak{}.\allowbreak{}.\allowbreak{})}
into the snapshot, for policies with JSON values. That works, but is hardly readable. A
custom \code{IAnnotation\allowbreak{}Code\allowbreak{}Generator} translates them into the
Fluent API:

\begin{sourcecode}[Shop\allowbreak{}Context\allowbreak{}Model\allowbreak{}Snapshot.\allowbreak{}cs]{csharp}
// Before
b.HasAnnotation("Acme:Rls:Enabled", true);
b.HasAnnotation("Acme:Rls:Forced", true);
b.HasAnnotation("Acme:Rls:Policy:tenant_isolation", "{\"Name\":\"tenant_isolation\",...}");

// After
RlsEntityTypeBuilderExtensions.HasRowLevelSecurity(b, true);
RlsEntityTypeBuilderExtensions.HasPolicy(b, "tenant_isolation", p => p
    .Using("\"TenantId\" = nullif(current_setting('app.tenant_id', true), '')::uuid")
    .WithCheck("\"TenantId\" = nullif(current_setting('app.tenant_id', true), '')::uuid"));
\end{sourcecode}

Since EF Core~7, EF writes extension methods in the snapshot as static calls
(\code{Extensions.\allowbreak{}Method(\allowbreak{}b,\allowbreak{}
.\allowbreak{}.\allowbreak{}.\allowbreak{})}) instead of
\code{b.\allowbreak{}Method(\allowbreak{}.\allowbreak{}.\allowbreak{}.\allowbreak{})}. This
avoids ambiguities when several providers or plugins define methods with the same name.
Tests that check the generated code should expect this form.

\begin{sourcecode}[Design/\allowbreak{}Internal/\allowbreak{}Tenancy\allowbreak{}Annotation\allowbreak{}Code\allowbreak{}Generator.\allowbreak{}cs (\allowbreak{}excerpt, simplified)]{csharp}
public override IReadOnlyList<MethodCallCodeFragment> GenerateFluentApiCalls(
    IEntityType entityType, IDictionary<string, IAnnotation> annotations)
{
    var calls = new List<MethodCallCodeFragment>();

    // CONTRACT: REMOVE translated annotations from the dictionary
    if (annotations.TryGetValue(RlsAnnotationNames.Enabled, out var enabled)
        && enabled.Value is true)
    {
        annotations.Remove(RlsAnnotationNames.Enabled);
        var forced = annotations.Remove(RlsAnnotationNames.Forced, out var f)
                     && f.Value is true;
        calls.Add(new MethodCallCodeFragment(HasRowLevelSecurityMethod, forced));
    }

    foreach (var key in annotations.Keys
                 .Where(k => k.StartsWith(RlsAnnotationNames.PolicyPrefix,
                                          StringComparison.Ordinal))
                 .ToList())                         // copy: the dictionary changes
    {
        var policy = RlsPolicy.Deserialize((string)annotations[key].Value!);
        annotations.Remove(key);
        calls.Add(new MethodCallCodeFragment(HasPolicyMethod, policy.Name,
            new NestedClosureCodeFragment("p", BuildPolicyChain(policy))));
    }

    calls.AddRange(base.GenerateFluentApiCalls(entityType, annotations));
    return calls;
}
\end{sourcecode}

\begin{datatable}{H{0.70}H{1.30}}
  \textbf{Pitfall} & \textbf{Solution} \\
  \tablehead
  Translated annotation stays in the dictionary & Appears twice in the snapshot. Always
  remove it; whatever remains, EF writes as \code{Has\allowbreak{}Annotation}. \\
  \tablerule
  Generic method in the \code{Method\allowbreak{}Info} & The snapshot uses the non-generic
  \code{Entity\allowbreak{}Type\allowbreak{}Builder}. Pick the non-generic overload via
  \code{Get\allowbreak{}Runtime\allowbreak{}Method} with exact parameter types. \\
  \tablerule
  \code{params} argument as an array &
  \code{Method\allowbreak{}Call\allowbreak{}Code\allowbreak{}Fragment} takes \code{params
  object?[]}. A \code{string[]} would become the argument list through array covariance:
  wrap it in \code{new object?[] \{ roles \}}. \\
  \tablerule
  Enums in the snapshot & Rendered as \code{Policy\allowbreak{}Command.\allowbreak{}Select}.
  The enum therefore lives in the always-imported namespace
  \code{Microsoft.\allowbreak{}Entity\allowbreak{}Framework\allowbreak{}Core}. \\
  \tablerule
  Lambda arguments &
  \code{Nested\allowbreak{}Closure\allowbreak{}Code\allowbreak{}Fragment(\allowbreak{}\textquotedbl{}p\textquotedbl{},\allowbreak{}
  chain)} produces \code{p =>
  p.\allowbreak{}Using(\allowbreak{}.\allowbreak{}.\allowbreak{}.\allowbreak{})}; chains are
  built with
  \code{fragment.\allowbreak{}Chain(\allowbreak{}.\allowbreak{}.\allowbreak{}.\allowbreak{})}.
  \\
  \tablerule
  Policy with default values only & Without expressions, roles and command the chain would be
  empty; \code{p => } without a call does not compile. \code{Build\allowbreak{}Policy\allowbreak{}Chain}
  then writes the default explicitly, \code{p => p.\allowbreak{}For(\allowbreak{}Policy\allowbreak{}Command.\allowbreak{}All)}:
  the same annotation, so no phantom migration. \\
\end{datatable}

The generator affects the snapshot, the designer files and the scaffolded context. The actual migrations keep
\code{.\allowbreak{}Annotation(\allowbreak{}.\allowbreak{}.\allowbreak{}.\allowbreak{})} on
\code{Alter\allowbreak{}Table}; that is the intended mechanism for changes.

% ------------------------------------------------------------
\chapter{Compiled models, trimming and NativeAOT}\label{ch:compiled-models-trimming-and-nativeaot}
\versionrisk{medium}

\code{dotnet ef dbcontext optimize} produces a \textbf{compiled model}: C\# code that builds
the runtime model directly instead of computing it at startup through conventions. This
shortens startup for large models considerably and is a prerequisite for NativeAOT. A plugin
must be prepared for it in four places:

\begin{datatable}{H{0.66}H{1.32}H{1.02}}
  \textbf{Building block} & \textbf{Requirement} & \textbf{In the sample} \\
  \tablehead
  Type mappings & Public type with a public static property \code{Default} (used for
  NativeAOT); converters/comparers as shared instances &
  \code{Tenant\allowbreak{}Id\allowbreak{}Type\allowbreak{}Mapping.\allowbreak{}Default},
  \code{Tenant\allowbreak{}Id\allowbreak{}Converter.\allowbreak{}Instance},
  \code{Tenant\allowbreak{}Id\allowbreak{}Value\allowbreak{}Comparer} \\
  \tablerule
  JSON reader/writer & Static \code{Instance} and \code{Constructor\allowbreak{}Expression}
  & \code{Tenant\allowbreak{}Id\allowbreak{}Json\allowbreak{}Reader\allowbreak{}Writer} \\
  \tablerule
  Value generators & Configured with
  \code{Has\allowbreak{}Value\allowbreak{}Generator\allowbreak{}Factory(\allowbreak{}Type)};
  factory with a public parameterless constructor &
  \code{Tenant\allowbreak{}Id\allowbreak{}Value\allowbreak{}Generator\allowbreak{}Factory}
  (chapter~\ref{ch:value-generation}) \\
  \tablerule
  Annotations & Simple values are taken over by EF automatically; an
  \code{ICSharp\allowbreak{}Runtime\allowbreak{}Annotation\allowbreak{}Code\allowbreak{}Generator}
  filters design-time annotations and translates complex values & Filters the RLS
  annotations, keeps the runtime annotation (chapter~\ref{ch:property-level-fluent-api-and}) \\
\end{datatable}

The last point is often misunderstood. EF takes annotations with simple values (strings,
numbers, bool, enums) into the compiled model \textbf{automatically}; the source of
\code{CSharp\allowbreak{}Runtime\allowbreak{}Annotation\allowbreak{}Code\allowbreak{}Generator}
shows this, writing them via \code{Generate\allowbreak{}Simple\allowbreak{}Annotations}. A
custom generator is therefore not needed to \emph{keep} annotations, but to \emph{filter
out} pure design-time annotations that nobody reads at run time, and to translate
annotations with complex values into code. Chapter~\ref{ch:property-level-fluent-api-and}
shows the filtering in the sample; the policies are stored there as JSON strings, that is,
as simple values.

\section{Trimming and AOT}

All packages of the plugin set
\code{<\allowbreak{}Is\allowbreak{}Aot\allowbreak{}Compatible>true<\allowbreak{}/\allowbreak{}Is\allowbreak{}Aot\allowbreak{}Compatible>}.
The compiler then reports code that uses reflection without annotations. The JSON
serialization of the policies uses a source generator instead of reflection:

% source: src/Acme.EntityFrameworkCore.Tenancy.PostgreSQL/Metadata/RlsPolicy.cs
\begin{sourcecode}[Metadata/\allowbreak{}Rls\allowbreak{}Policy.\allowbreak{}cs (\allowbreak{}excerpt)]{csharp}
[JsonSourceGenerationOptions(UseStringEnumConverter = true)]
[JsonSerializable(typeof(RlsPolicy))]
internal partial class RlsJsonContext : JsonSerializerContext;
\end{sourcecode}

Where dynamic code is unavoidable, the warning is suppressed with a justification instead of
being ignored. The isolation convention builds the query filter as an expression tree, which
needs dynamic code. Conventions only run while the model is built, however, and NativeAOT
requires a compiled model, in which they no longer run; the suppression states exactly this
reason. The warnings only show reliably in a full build; incremental builds keep quiet about
them (upgrade playbook in chapter~\ref{ch:versioning-and-compatibility}).

\begin{note}
  Conventions, snapshots and migrations only run at design time or during model building;
  with a compiled model and precompiled queries they disappear at run time. For AOT, the
  critical parts are therefore type mappings, converters and interceptors; with precompiled
  queries, translators already run at build time.
\end{note}

\section{What is actually checked as a consumer}\label{sec:consumer-test}

\code{Is\allowbreak{}Aot\allowbreak{}Compatible} runs analyzers but proves nothing at run time.
A unit test that generates compiled model code and checks text in it
(\code{Compiled\allowbreak{}Model\allowbreak{}Tests}) does not compile or load the model either.
\code{eng/\allowbreak{}consumer-\allowbreak{}test/\allowbreak{}run.\allowbreak{}sh} therefore
checks the plugin the way an application gets it, in a temporary directory that is deleted at
the end:

\begin{datatable}{H{0.55}H{1.45}}
  \textbf{Step} & \textbf{Check} \\
  \tablehead
  Packages & \code{dotnet pack} of the four libraries (core, relational layer,
  PostgreSQL, SQL Server) into a local source; a fresh project
  references the packages, not the projects. Its own NuGet cache, package source mapping with
  \code{Acme.*} only from the local source and a SHA-512 comparison prove that exactly the
  packages just built are resolved. \\
  \tablerule
  Tooling & \code{dotnet-ef} as a local tool. The design-time services are loaded by the
  attribute in the startup project. SQL script and snapshot contain RLS, \code{FORCE}, policy and index (PostgreSQL) and
  security policy and session key (SQL Server).
  \code{has-\allowbreak{}pending-\allowbreak{}model-\allowbreak{}changes} then reports no
  change. \\
  \tablerule
  Runtime model & \code{TenantId} becomes \code{uuid}, the runtime annotation names the
  property, \code{Current\allowbreak{}Tenant(\allowbreak{})} and the query filter are
  translated. \\
  \tablerule
  Compiled model & \code{dotnet ef dbcontext optimize}, compile, load with
  \code{Use\allowbreak{}Model}, the same checks. \\
  \tablerule
  NativeAOT (\code{-{}-nativeaot}) & Publish with \code{Publish\allowbreak{}Aot} and
  precompiled queries, then run against PostgreSQL in Docker: two tenants write through the
  interceptors, one reads, RLS returns only its rows. No trim or AOT warning may come from an
  Acme assembly. The deployment sequence from chapter~\ref{ch:retrofitting-and-operating} runs
  as well: roles, migration bundle as the owner, no DDL for the app role. \\
\end{datatable}

CI runs the flow in a job of its own with the newest stable 10.0.x patch for the tooling and
prints the versions. The test found things the unit tests could not show:

\begin{itemize}
  \item \textbf{Query filters and compiled models exclude each other.} EF Core~10 rejects
    \code{optimize}: “Compiled model can't be generated, because query filters are not
    supported.” Whoever enables the plugin's optional filter gives up compiled model and
    NativeAOT; protection then comes from RLS alone. The test expects this error explicitly, so
    that a later EF version that supports filters stands out.
  \item \textbf{Value comparers through reflection.} The type mapping for \code{TenantId} brought
    no comparers of its own. The compiled model then only refers to
    \code{Tenant\allowbreak{}Id\allowbreak{}Type\allowbreak{}Mapping.\allowbreak{}Default}, and
    EF creates the comparer on first tracking with \code{Make\allowbreak{}Generic\allowbreak{}Method},
    which made the NativeAOT application fail. Since then both mappings carry shared instances
    from \code{Tenant\allowbreak{}Id\allowbreak{}Value\allowbreak{}Comparer}, like converter and
    JSON reader.
  \item \textbf{Limits of EF itself.} Precompiled queries and NativeAOT are experimental in EF
    Core~10 (warning EF9100). \code{Publish\allowbreak{}Aot} must not be in the project, because
    EF otherwise refuses to build the model at design time already. The package
    \code{Microsoft.\allowbreak{}Entity\allowbreak{}Framework\allowbreak{}Core.\allowbreak{}Tasks}
    generates model and queries on publish; an additional, manually generated model gets in its
    way. Precompilation rejects \code{o.\allowbreak{}Total > 0} on a \code{decimal} property,
    \code{0m} works. EF, EF Relational and Npgsql report trim and AOT warnings themselves
    (IL2104, IL3053). With EF Design~10.0.0, publishing fails because the Tasks package treats
    the NuGet audit warnings as errors; that is why the test uses the newest patch by default.
\end{itemize}

Not checked are SQL Server under NativeAOT (the package is marked for it, a run with a real
database is missing) and the Cosmos package, which the consumer test does not pack.


% ------------------------------------------------------------
\chapter{Scaffolding (reverse engineering)}\label{ch:scaffolding-reverse-engineering}
\versionrisk{medium}

\code{dotnet ef dbcontext scaffold} reads an existing database schema and generates entities
and a context. A plugin that manages its own schema objects can recognize them during this
too. The route goes through two design-time services:

\begin{datatable}{H{0.60}H{1.40}}
  \textbf{Service} & \textbf{Task} \\
  \tablehead
  \code{IDatabase\allowbreak{}Model\allowbreak{}Factory} & Reads the schema into a
  \code{Database\allowbreak{}Model}. A derived Npgsql factory queries
  \code{pg\_\allowbreak{}policies} and
  \code{pg\_\allowbreak{}class.\allowbreak{}relrowsecurity} and attaches them as annotations
  to \code{Database\allowbreak{}Table}. \\
  \tablerule
  \code{IAnnotation\allowbreak{}Code\allowbreak{}Generator} & Writes these annotations as
  Fluent API into the generated \code{On\allowbreak{}Model\allowbreak{}Creating}, with the
  same logic as for the snapshot. \\
\end{datatable}

The sample implements both: a custom database model factory reads RLS state and policies
from the database, and the annotation code generator from
chapter~\ref{ch:design-time-services} writes them into the scaffolded context.
Chapter~\ref{ch:scaffolding-in-practice} shows the implementation in detail.
